\documentclass[10pt]{article}
% Shared LaTeX style for MIT 18.06 exam revision notes.
% Compile topic files with Tectonic, XeLaTeX, or LuaLaTeX.

\usepackage[a4paper,margin=0.72in]{geometry}
\usepackage{fontspec}
\usepackage{amsmath,amssymb,mathtools}
\usepackage{booktabs,array,tabularx}
\usepackage{enumitem}
\usepackage{titlesec}
\usepackage{fancyhdr}
\usepackage{microtype}
\usepackage{xcolor}

\setmainfont{Times New Roman}
\setsansfont{Arial}

\setlength{\parindent}{0pt}
\setlength{\parskip}{3.6pt}
\setlength{\abovedisplayskip}{6pt}
\setlength{\belowdisplayskip}{6pt}
\setlength{\abovedisplayshortskip}{4pt}
\setlength{\belowdisplayshortskip}{4pt}
\renewcommand{\arraystretch}{1.08}
\setlength{\tabcolsep}{4.5pt}

\titleformat{\section}
  {\normalfont\large\bfseries}
  {\thesection}
  {0.55em}
  {}
\titlespacing*{\section}{0pt}{10pt}{3pt}

\titleformat{\subsection}
  {\normalfont\normalsize\bfseries}
  {\thesubsection}
  {0.5em}
  {}
\titlespacing*{\subsection}{0pt}{7pt}{2pt}

\setlist[itemize]{leftmargin=1.35em,itemsep=1pt,topsep=2pt,parsep=0pt}
\setlist[enumerate]{leftmargin=1.55em,itemsep=1pt,topsep=2pt,parsep=0pt}

\pagestyle{fancy}
\fancyhf{}
\fancyfoot[C]{\scriptsize\color{black!45}MIT 18.06 Linear Algebra Exam Notes --- \thepage}
\renewcommand{\headrulewidth}{0pt}
\renewcommand{\footrulewidth}{0pt}

\newcommand{\notesubtitle}[1]{%
  {\centering\itshape #1\par}
  \vspace{2pt}
}

\newcommand{\source}[1]{%
  {\centering\small #1\par}
  \vspace{6pt}
}

\newcommand{\labelnote}[2]{%
  \par\smallskip
  \noindent\textbf{#1.}\quad #2
  \par\smallskip
}

\newcommand{\tightlabel}[1]{%
  \par\smallskip
  \noindent\textbf{#1.}\quad
}

\newcolumntype{Y}{>{\raggedright\arraybackslash}X}
\newcolumntype{L}[1]{>{\raggedright\arraybackslash}p{#1}}

\newenvironment{notetable}
  {\small\begin{tabularx}{\linewidth}}
  {\end{tabularx}}


\begin{document}

\begin{center}
{\LARGE 4.1 Orthogonality and the Four Fundamental Subspaces\par}
\end{center}
\notesubtitle{Concise exam revision notes for orthogonal complements in \(\mathbb{R}^n\) and \(\mathbb{R}^m\).}
\source{Based on handwritten MIT 18.06 revision notes.}

\section{Core Idea}

\labelnote{Core idea}{Orthogonality connects the four fundamental subspaces in pairs. Row space is perpendicular to nullspace; column space is perpendicular to left nullspace.}

The dot product test is
\[
v^{T}w=0.
\]
If this holds for every vector \(v\) in one subspace and every vector \(w\) in another, the two subspaces are orthogonal.

\section{Orthogonal Subspaces vs Orthogonal Complements}

Orthogonal subspaces may intersect only at zero, but that alone does not mean they fill the whole space.

\begin{center}
\small
\begin{tabularx}{\linewidth}{@{}L{0.30\linewidth}Y Y@{}}
\toprule
Idea & Meaning & Exam check \\
\midrule
Orthogonal subspaces & Every vector in one is perpendicular to every vector in the other. & Dot products are zero. \\
Orthogonal complements & Orthogonal subspaces whose dimensions add to the ambient dimension. & Dot products zero and dimensions add correctly. \\
\bottomrule
\end{tabularx}
\end{center}

\labelnote{Common mistake}{Do not say ``orthogonal complement'' from perpendicularity alone. Also check the dimension count.}

\section{Row Space and Nullspace}

Let \(A\in\mathbb{R}^{m\times n}\). Each row of \(A\) is a vector in \(\mathbb{R}^n\). If \(x\in N(A)\), then
\[
Ax=0.
\]
This means every row of \(A\) has dot product zero with \(x\). Therefore
\[
C(A^{T})\perp N(A).
\]

Their dimensions add to the whole ambient space \(\mathbb{R}^n\):
\[
\dim C(A^{T})+\dim N(A)=r+(n-r)=n.
\]
So they are orthogonal complements:
\[
N(A)=C(A^{T})^{\perp},
\qquad
C(A^{T})=N(A)^{\perp}.
\]

\section{Column Space and Left Nullspace}

The left nullspace is
\[
N(A^{T})=\{y:A^{T}y=0\}.
\]
If \(y\in N(A^{T})\), then \(y\) is perpendicular to every column of \(A\). Therefore
\[
C(A)\perp N(A^{T}).
\]

Their dimensions add to the whole ambient space \(\mathbb{R}^m\):
\[
\dim C(A)+\dim N(A^{T})=r+(m-r)=m.
\]
So they are orthogonal complements:
\[
N(A^{T})=C(A)^{\perp},
\qquad
C(A)=N(A^{T})^{\perp}.
\]

\section{Four-Subspace Summary}

\begin{center}
\small
\begin{tabularx}{\linewidth}{@{}L{0.25\linewidth}L{0.16\linewidth}L{0.18\linewidth}Y@{}}
\toprule
Subspace & Lives in & Dimension & Orthogonal complement \\
\midrule
Row space \(C(A^{T})\) & \(\mathbb{R}^n\) & \(r\) & Nullspace \(N(A)\). \\
Nullspace \(N(A)\) & \(\mathbb{R}^n\) & \(n-r\) & Row space \(C(A^{T})\). \\
Column space \(C(A)\) & \(\mathbb{R}^m\) & \(r\) & Left nullspace \(N(A^{T})\). \\
Left nullspace \(N(A^{T})\) & \(\mathbb{R}^m\) & \(m-r\) & Column space \(C(A)\). \\
\bottomrule
\end{tabularx}
\end{center}

\labelnote{Exam memory}{The complement must live in the same ambient space: row with null in \(\mathbb{R}^n\), column with left null in \(\mathbb{R}^m\).}

\section{Small Example}

Let
\[
A=
\begin{bmatrix}
1&2&0\\
0&0&1
\end{bmatrix}.
\]
The row space is
\[
C(A^{T})=\operatorname{span}\{(1,2,0),(0,0,1)\}.
\]
The nullspace is found from
\[
x_1+2x_2=0,
\qquad
x_3=0,
\]
so
\[
N(A)=\operatorname{span}\{(-2,1,0)\}.
\]
The dot products are zero:
\[
(1,2,0)\cdot(-2,1,0)=0,
\qquad
(0,0,1)\cdot(-2,1,0)=0.
\]
This shows the row space direction and nullspace direction are perpendicular in \(\mathbb{R}^3\).

\section{How to Use This in Problems}

\begin{center}
\small
\begin{tabularx}{\linewidth}{@{}L{0.32\linewidth}Y@{}}
\toprule
If you know a vector is perpendicular to & Then place it in \\
\midrule
Every row of \(A\) & \(N(A)\). \\
Every vector in \(N(A)\) & \(C(A^{T})\). \\
Every column of \(A\) & \(N(A^{T})\). \\
Every vector in \(N(A^{T})\) & \(C(A)\). \\
\bottomrule
\end{tabularx}
\end{center}

\labelnote{Dimension check}{Perpendicularity plus the correct dimension identifies the whole orthogonal complement.}

\section{Common Mistakes}

\begin{center}
\small
\begin{tabularx}{\linewidth}{@{}L{0.40\linewidth}Y@{}}
\toprule
Mistake & Correct exam habit \\
\midrule
Putting row space in \(\mathbb{R}^m\). & Row vectors have length \(n\), so row space lives in \(\mathbb{R}^n\). \\
Pairing column space with \(N(A)\). & Column space pairs with \(N(A^{T})\). \\
Forgetting the left nullspace. & Solve \(A^{T}y=0\), not \(Ax=0\). \\
Calling any perpendicular subspaces complements. & Check that the dimensions add to \(n\) or \(m\). \\
Copying the picture instead of using the dimension count. & Use the table and formulas to reconstruct the relationship. \\
\bottomrule
\end{tabularx}
\end{center}

\section{Final Exam Checklist}

\tightlabel{Checklist}
\begin{enumerate}
  \item Write the size: \(A\) is \(m\times n\).
  \item Identify the ambient space: \(\mathbb{R}^n\) or \(\mathbb{R}^m\).
  \item Pair \(C(A^{T})\) with \(N(A)\) inside \(\mathbb{R}^n\).
  \item Pair \(C(A)\) with \(N(A^{T})\) inside \(\mathbb{R}^m\).
  \item Use \(Ax=0\) for the nullspace.
  \item Use \(A^{T}y=0\) for the left nullspace.
  \item Check dimension sums: \(r+(n-r)=n\) and \(r+(m-r)=m\).
  \item State the orthogonal complement relationship directly.
\end{enumerate}

\end{document}
